\section{Conclusion}

Stablecoins transform parallel foreign exchange markets not only by lowering transaction costs, but by fundamentally altering how information is aggregated and acted upon. By generating a high-frequency public signal, they compress belief dispersion and enable coordinated behavior among agents.

This creates a fundamental tension between efficiency and stability. While stablecoins improve access to foreign currency and enhance price discovery, they can also increase the likelihood of coordinated exits in fragile macroeconomic environments. 

More broadly, our results suggest that financial innovations that improve transparency may have unintended destabilizing effects when agents’ actions are strategically complementary. Understanding this tradeoff is essential for the design of policy responses to the growing role of digital financial instruments in emerging markets.